\chapter{Lesson 2: Colours and Clothing}

\section{Vocabulary}

\subsection{Colours}

Colour words are adjectives, but most can also take \zh{色} \emph{sè} ``colour'' and
act as nouns: \zh{红色}, \zh{蓝色}. \zh{咖啡色} ``coffee colour'' is brown.

\begin{center}
\begin{tikzpicture}[sw/.style={draw, minimum width=1.25cm, minimum height=0.7cm}]
  \foreach \i/\c/\zhc/\py/\en in {
      0/red!80!black/红/hóng/red,
      1/yellow!90!orange/黄/huáng/yellow,
      2/blue!70/蓝/lán/blue,
      3/green!55!black/绿/lǜ/green,
      4/white/白/bái/white,
      5/black/黑/hēi/black,
      6/brown!80!black/咖啡色/kāfēisè/brown,
      7/gray/灰/huī/grey,
      8/violet!80!black/紫/zǐ/purple,
      9/pink/粉红/fěnhóng/pink,
      10/orange/橙/chéng/orange} {
    \pgfmathsetmacro{\x}{mod(\i,6)*2.5}
    \pgfmathsetmacro{\y}{-floor(\i/6)*1.9}
    \node[sw, fill=\c] at (\x,\y) {};
    \node[font=\small] at (\x,\y-0.6) {\zh{\zhc} \emph{\py}};
    \node[font=\scriptsize] at (\x,\y-0.95) {\en};
  }
\end{tikzpicture}
\end{center}

\subsection{Clothing}

\begin{center}
\begin{tabular}{llll}
\toprule
Characters & Pinyin & Measure word & Meaning \\
\midrule
\zh{衣服} & yīfu & \zh{件} & clothes \\
\zh{衬衫} & chènshān & \zh{件} & shirt \\
\zh{T恤} & T-xù & \zh{件} & T-shirt \\
\zh{毛衣} & máoyī & \zh{件} & sweater \\
\zh{大衣} & dàyī & \zh{件} & overcoat \\
\zh{裤子} & kùzi & \zh{条} & pants \\
\zh{牛仔裤} & niúzǎikù & \zh{条} & jeans \\
\zh{裙子} & qúnzi & \zh{条} & skirt, dress \\
\zh{鞋} & xié & \zh{双} & shoes \\
\zh{袜子} & wàzi & \zh{双} & socks \\
\zh{帽子} & màozi & \zh{顶} dǐng & hat \\
\zh{眼镜} & yǎnjìng & \zh{副} fù & glasses \\
\zh{手套} & shǒutào & \zh{副}/\zh{双} & gloves \\
\bottomrule
\end{tabular}
\end{center}

\subsection{Other words}

\begin{center}
\begin{tabular}{llll}
\toprule
Characters & Pinyin & Part of speech & Meaning \\
\midrule
\zh{穿} & chuān & verb & to wear (clothes, shoes) \\
\zh{戴} & dài & verb & to wear (accessories) \\
\zh{换} & huàn & verb & to change, exchange \\
\zh{合适} & héshì & adjective & suitable, fitting \\
\zh{好看} & hǎokàn & adjective & good-looking \\
\zh{长} & cháng & adjective & long \\
\zh{短} & duǎn & adjective & short \\
\zh{号} & hào & noun & size \\
\zh{虽然} & suīrán & conjunction & although \\
\zh{但是} & dànshì & conjunction & but \\
\zh{又\ldots{}又\ldots} & yòu \ldots yòu \ldots & & both \ldots and \ldots \\
\zh{一样} & yíyàng & adjective & the same \\
\zh{不但\ldots{}而且\ldots} & búdàn \ldots érqiě \ldots & & not only \ldots but also \ldots \\
\zh{质量} & zhìliàng & noun & quality \\
\zh{样子} & yàngzi & noun & style, appearance \\
\bottomrule
\end{tabular}
\end{center}

\section{\zh{穿} vs \zh{戴}}

\begin{important}[Two verbs for ``wear'']
\zh{穿} is for things you get \emph{into}: clothes, pants, socks, shoes.
\zh{戴} is for things you \emph{put on} or attach: hats, glasses, watches, gloves,
jewellery, masks.
\begin{itemize}
  \item \zh{他穿着一件黑大衣，戴着一顶红帽子。} He's wearing a black coat and a red hat.
  \item \zh{冬天在滑铁卢得戴手套。} In winter in Waterloo you have to wear gloves.
\end{itemize}
\end{important}

\section{Grammar}

\subsection{\zh{又\ldots{}又\ldots}}

\zh{又} A \zh{又} B ``both A and B'' links two adjectives (or verb phrases) that point the
same way, both positive or both negative.
\begin{itemize}
  \item \zh{这条裤子又便宜又好看。} These pants are both cheap and nice.
  \item \zh{那双鞋又贵又不舒服。} Those shoes are expensive and uncomfortable.
  \item *\zh{又便宜又不好看} sounds odd: mixed evaluations need \zh{虽然\ldots{}但是}.
\end{itemize}

\subsection{\zh{虽然\ldots{}但是/可是\ldots}}

Unlike English, Chinese marks both halves of a concession.
\begin{itemize}
  \item \zh{虽然这件毛衣很贵，但是质量很好。} Although this sweater is expensive, the
        quality is good.
  \item \zh{她虽然很喜欢那条裙子，可是没买。} Although she likes that dress, she didn't
        buy it.
\end{itemize}
\zh{虽然} can come before or after the subject.

\subsection{\zh{跟\ldots{}一样}}

A \zh{跟} B \zh{一样} ``A is the same as B''; add an adjective for ``as \ldots as''.
\begin{itemize}
  \item \zh{我的衬衫跟你的一样。} My shirt is the same as yours.
  \item \zh{这件跟那件一样大。} This one is as big as that one.
  \item \zh{这条裤子跟那条不一样长。} These pants are not the same length as those.
\end{itemize}

\subsection{\zh{不但\ldots{}而且\ldots}}

\begin{itemize}
  \item \zh{这家商店的衣服不但便宜，而且样子很好看。} Not only are this store's clothes
        cheap, they're also stylish.
\end{itemize}
When the two clauses have different subjects, \zh{不但} goes before the first subject.

\subsection{Trying on: verb reduplication and \zh{一下}}

\zh{试试}, \zh{试一下}, \zh{看看}: ``give it a try / take a look''.
\begin{itemize}
  \item \zh{我可以试试这件吗？}
  \item \zh{太长了，能不能换一条短一点儿的？} It's too long; could I exchange it for a
        slightly shorter one?
\end{itemize}

\subsection{Sizes}

Chinese sizes: \zh{大号} L, \zh{中号} M, \zh{小号} S. Shoe sizes in China use European
numbers (\zh{四十二号}). \zh{你穿多大号的？} ``What size do you wear?''

\section{Colour symbolism}

\begin{aside}
\begin{itemize}
  \item \zh{红} stands for luck and celebration: weddings, New Year couplets, red
        envelopes (\zh{红包} hóngbāo) of money.
  \item \zh{白} is the colour of mourning; wearing all white to a wedding or giving white
        flowers is a faux pas.
  \item \zh{黄} was the imperial colour; today \zh{黄色} can also mean ``pornographic''
        (\zh{黄色电影}).
  \item \zh{戴绿帽子} ``to wear a green hat'' means one's partner is unfaithful, so no one
        buys green hats as gifts.
\end{itemize}
Traditional forms: \zht{紅} red, \zht{綠} green, \zht{顏色} colour.
\end{aside}

\section{Model conversation}

\begin{itemize}
  \item \textbf{A:} \zh{你看，这件咖啡色的大衣怎么样？}
  \item \textbf{B:} \zh{样子不错，可是有点儿长。你穿多大号的？}
  \item \textbf{A:} \zh{中号。我试一下。}
  \item \textbf{B:} \zh{虽然颜色很好看，但是大衣太长了。换一件短一点儿的吧。}
  \item \textbf{A:} \zh{这件跟那件一样，不过短一点儿。又便宜又合适！就买这件。}
\end{itemize}
(\zh{不过} búguò = however.)

\section*{Review}
\begin{itemize}
  \item Colours: \zh{红 黄 蓝 绿 白 黑 咖啡色 灰 紫 粉红 橙}.
  \item \zh{穿} clothes and shoes; \zh{戴} accessories.
  \item Clothing measure words: \zh{件} tops, \zh{条} pants/skirts, \zh{双} pairs.
  \item \zh{又}A\zh{又}B (same direction); \zh{虽然\ldots{}但是} (contrast).
  \item A \zh{跟} B \zh{一样}(+ adj); \zh{不但\ldots{}而且}.
\end{itemize}
